%% TinyPrognostics — Paper Skeleton
%% Journal: Artificial Intelligence for Engineering (Wiley/IET)
%% Template: WileyNJDv5 | AMA style | Times font | 1-column
%% Compile with: XeLaTeX (preferred) or pdfLaTeX
%%
%% ============================================================
%% IMPORTANT: Download WileyNJDv5.cls and place it in the same
%% folder as this file. Get it from:
%% https://authorservices.wiley.com/author-resources/Journal-Authors/Prepare/latex.html
%% ============================================================

\documentclass[AMA,Times1COL]{WileyNJDv5}

%% ---------- Extra packages ----------
\usepackage{amsmath}
\usepackage{amssymb}
\usepackage{graphicx}
\usepackage{booktabs}
\usepackage{multirow}
\usepackage{algorithm}
\usepackage{algpseudocode}
\usepackage{url}
\usepackage{hyperref}
\hypersetup{colorlinks=true, linkcolor=blue, citecolor=blue, urlcolor=blue}

%% ---------- Journal metadata ----------
\articletype{ORIGINAL RESEARCH}

\received{00 Month 2026}
\revised{00 Month 2026}
\accepted{00 Month 2026}

\raggedbottom

%% ============================================================
\begin{document}
%% ============================================================

\title{NanoSentry: An Edge-Deployable Unified Neural Architecture
for Multi-Domain Industrial Prognostics and Health Management}

%% --- Authors ---
\author[1]{MD Hamid Borkot Tulla}

%% --- Affiliations ---
\authormark{TULLA}

\address[1]{\orgdiv{Department Name},
\orgname{Institution Name},
\orgaddress{\state{State}, \country{Country}}}

%% --- Correspondence ---
\corres{MD Hamid Borkot Tulla, \email{your@email.com}}

%% --- Abstract ---
\abstract{
Prognostics and health management (PHM) systems deployed on edge hardware face
a fundamental conflict: state-of-the-art deep learning models for remaining
useful life (RUL) estimation and fault classification demand computational
resources that microcontroller-class devices cannot provide. We introduce
NanoSentry, a 12{,}052-parameter (47.1\,KB) unified neural architecture that
resolves this conflict by jointly performing RUL regression and health-state
classification across heterogeneous industrial domains within a single model.
NanoSentry combines dilated causal 1D convolutions with a novel CrossSensorGate
mechanism for adaptive inter-sensor feature weighting and a GRU with orthogonal
initialisation for temporal context. Evaluated on three public benchmarks ---
NASA C-MAPSS turbofan engine degradation (four subsets), NASA lithium-ion
battery state-of-health, and CWRU rolling bearing fault classification ---
NanoSentry achieves RMSE of 14.75 on C-MAPSS FD001, surpassing an LSTM
baseline 4.5$\times$ its size, and 99.96\% fault classification accuracy.
Critically, in the operationally decisive window of RUL $\leq$ 30 cycles,
NanoSentry reduces prediction error by 3--7$\times$ compared with a
comparable-size CNN baseline whose critical-zone RMSE pathologically
exceeds its overall error --- a hazardous property for safety-critical
maintenance. To our knowledge, NanoSentry is the first sub-50\,KB architecture
validated simultaneously across turbofan, battery, and bearing prognostic tasks.
}

%% --- Keywords ---
\keywords{prognostics and health management, remaining useful life estimation,
edge deployment, dilated temporal convolution, CrossSensorGate, multi-task
learning, bearing fault diagnosis, battery state-of-health}

\maketitle

%% ============================================================
\section{Introduction}\label{sec:intro}
%% ============================================================

\subsection{Industrial Motivation}\label{sec:intro:motivation}
% TODO: Write 3--4 sentences on cost of unplanned failure, IIoT constraints,
% and the push toward on-device PHM inference.
% Cite: Lei et al.~\cite{Lei2018}, Zhang et al.~\cite{Zhang2019}

\subsection{The Model-Size Problem}\label{sec:intro:size}
% TODO: Explain why LSTM/Transformer models cannot run on MCU-class hardware.
% Cite: Zheng et al.~\cite{Zheng2017}, Mo et al.~\cite{Mo2021}

\subsection{The Multi-Domain Gap}\label{sec:intro:gap}
% TODO: No single sub-50\,KB model shown on turbofan + battery + bearing.
% Reference the Wu lightweight CNN paper as partial prior work.

\subsection{Contributions}\label{sec:intro:contributions}
The main contributions of this work are as follows:
\begin{enumerate}
 \item We propose NanoSentry, a 12{,}052-parameter unified PHM architecture
 that jointly performs RUL regression and health-state classification across
 three industrial domains within a single 47.1\,KB model.
 \item We introduce the CrossSensorGate mechanism, a lightweight attention
 module that adaptively weights inter-sensor features without increasing
 model depth or parameter count.
 \item We demonstrate, for the first time, that a sub-50\,KB architecture
 achieves competitive accuracy against models 4.5$\times$ its size on the
 NASA C-MAPSS benchmark.
 \item We establish a critical-zone RMSE evaluation protocol (RUL $\leq$ 30
 cycles) and show that NanoSentry reduces end-of-life prediction error by
 3--7$\times$ compared with a comparable-size CNN baseline.
\end{enumerate}

%% ============================================================
\section{Related Work}\label{sec:related}
%% ============================================================

\subsection{RUL Estimation for Turbofan Engines}\label{sec:related:rul}
% TODO: Summarise CNN, LSTM, TCN, Transformer approaches on C-MAPSS.
% Cite: Zheng2017, Elsheikh2019, Li2020, Zhang2023atcn, Mo2021,
% Chen2024, Zhang2023imdssn, Lin2023, Liu2021, Gong2022

\subsection{Battery State-of-Health Estimation}\label{sec:related:battery}
% TODO: Data-driven SOH methods.
% Cite: Wang2024, Gong2024, Richardson2017

\subsection{Bearing Fault Diagnosis}\label{sec:related:bearing}
% TODO: CNN and deep learning fault classification.
% Cite: Cheng2020, Han2021, Zhao2020, Wen2018, Zhang2018, Liao2023,
% Liang2022

\subsection{Multi-Task and Unified PHM}\label{sec:related:multitask}
% TODO: Joint RUL + fault models.
% Cite: Zhou2025, Gong2022

\subsection{Lightweight Neural Architectures for PHM}\label{sec:related:efficient}
% TODO: Existing compact models, their limitations.
% End with the gap statement:
% ``No published work has demonstrated a unified architecture below 50\,KB
% validated simultaneously on turbofan, battery, and bearing PHM tasks.''

%% ============================================================
\section{NanoSentry Architecture}\label{sec:arch}
%% ============================================================

\subsection{Problem Formulation}\label{sec:arch:problem}
% Given a multivariate time-series window X of length T and C sensors,
% predict RUL scalar and health-state class simultaneously.
% TODO: Write formal problem statement.

\subsection{Architecture Overview}\label{sec:arch:overview}
% TODO: Reference Figure~\ref{fig:arch} and describe the four-stage pipeline:
% Linear embed -> TCN stack -> CrossSensorGate + Skip -> GRU -> dual heads.

% ----- FIGURE 1: Architecture diagram -----
\begin{figure}[!t]
 \centering
 % \includegraphics[width=\columnwidth]{figs/architecture.pdf}
 \fbox{\parbox{0.9\columnwidth}{\centering\vspace{2cm}
 [PLACEHOLDER: NanoSentry architecture diagram]\vspace{2cm}}}
 \caption{NanoSentry architecture. Input windows pass through a linear
 embedding, three dilated causal convolution blocks (dilation 1, 2, 4),
 a CrossSensorGate with residual skip, and a GRU before branching into
 the RUL regression head and health-state classification head.}
 \label{fig:arch}
\end{figure}

\subsection{Dilated Causal Convolution Block}\label{sec:arch:dilconv}
% TODO: Describe the TCN block.
Let $\mathbf{X} \in \mathbb{R}^{T \times C}$ denote an input window of
$T$ timesteps and $C$ sensor channels. A dilated causal convolution with
dilation factor $d$ and kernel size $k$ computes
\begin{equation}\label{eq:dilconv}
 \mathbf{h}_t = \text{ReLU}\!\left(
 \sum_{j=0}^{k-1} \mathbf{W}_j \, \mathbf{x}_{t - d \cdot j}
 + \mathbf{b} \right),
\end{equation}
where $\mathbf{W}_j$ are learnable filter weights and causality is enforced
by restricting attention to $t - d\cdot j \leq t$. Three successive blocks
with $d \in \{1,2,4\}$ yield a receptive field of $7(k-1)+1$ timesteps
without pooling.

\subsection{CrossSensorGate}\label{sec:arch:gate}
% TODO: Describe and motivate the gate.
Let $\mathbf{H} \in \mathbb{R}^{T \times d_{\text{model}}}$ be the TCN
output. The CrossSensorGate computes a channel-wise attention mask
\begin{equation}\label{eq:gate}
 \mathbf{g} = \sigma\!\left( \mathbf{W}_g \mathbf{H}^\top + \mathbf{b}_g \right),
 \qquad
 \hat{\mathbf{H}} = \mathbf{g}^\top \odot \mathbf{H},
\end{equation}
where $\sigma$ is the sigmoid function, $\odot$ denotes element-wise
multiplication, and $\mathbf{W}_g \in \mathbb{R}^{d_{\text{model}}
\times d_{\text{model}}}$. The gate suppresses sensor channels
irrelevant to the current degradation regime without adding layers.

\subsection{Residual Skip Connection}\label{sec:arch:skip}
A linear projection $\mathbf{W}_{\text{skip}} \in
\mathbb{R}^{C \times d_{\text{model}}}$ maps the raw input directly to
the model dimension, and the gated feature map is augmented by
\begin{equation}\label{eq:skip}
 \tilde{\mathbf{H}} = \hat{\mathbf{H}} + \mathbf{X}\mathbf{W}_{\text{skip}}.
\end{equation}
This preserves raw sensor gradients during backpropagation and
prevents vanishing-gradient degradation in shallow stacks.

\subsection{Gated Recurrent Unit}\label{sec:arch:gru}
The skip-augmented representation is passed to a single-layer GRU
\begin{equation}\label{eq:gru}
 \mathbf{z}_t = \text{GRU}\!\left(\tilde{\mathbf{h}}_t,\, \mathbf{z}_{t-1}\right),
\end{equation}
with hidden dimension $d = 24$ and weights initialised with an
orthogonal matrix~\cite{TODO} to prevent gradient explosion at the
start of training.

\subsection{Dual Output Heads}\label{sec:arch:heads}
The final hidden state $\mathbf{z}_T$ is fed to two independent
fully-connected heads:
\begin{equation}\label{eq:heads}
 \hat{y}_{\text{RUL}} = \mathbf{w}_r^\top \mathbf{z}_T + b_r,
 \qquad
 \hat{\mathbf{y}}_{\text{state}} = \text{softmax}\!\left(
 \mathbf{W}_c \mathbf{z}_T + \mathbf{b}_c \right).
\end{equation}

\subsection{Joint Loss Function}\label{sec:arch:loss}
The model is trained end-to-end by minimising
\begin{equation}\label{eq:loss}
 \mathcal{L} = \lambda_1 \, \mathcal{L}_{\text{RUL}}
 + \lambda_2 \, \mathcal{L}_{\text{state}},
\end{equation}
where $\mathcal{L}_{\text{RUL}}$ is mean squared error,
$\mathcal{L}_{\text{state}}$ is cross-entropy, and
$\lambda_1, \lambda_2$ are task weights set to 1.0 and 0.5 respectively.
The asymmetric NASA scoring function
\begin{equation}\label{eq:nasa}
 s_i =
\begin{cases}
 e^{-d_i / 13} - 1 & d_i < 0 \quad \text{(early prediction)}\\
 e^{\phantom{-}d_i / 10} - 1 & d_i \geq 0 \quad \text{(late prediction)}
\end{cases}
\end{equation}
is reported as a secondary metric to penalise late predictions more
severely, where $d_i = \hat{y}_i - y_i$.

\subsection{Training Procedure}\label{sec:arch:training}
\begin{algorithm}[!h]
\caption{NanoSentry training procedure}\label{alg:training}
\begin{algorithmic}[1]
\Require Dataset $\mathcal{D}$, epochs $E$, batch size $B$,
 learning rate $\eta$, task weights $\lambda_1, \lambda_2$
\Ensure Trained model parameters $\boldsymbol{\theta}$
\State Initialise $\boldsymbol{\theta}$ (GRU weights: orthogonal)
\State Standardise inputs per sensor channel
\For{epoch $e = 1, \ldots, E$}
 \For{each mini-batch $(\mathbf{X}, y_{\text{RUL}}, y_{\text{state}}) \in \mathcal{D}$}
 \State Forward pass: compute $\hat{y}_{\text{RUL}}$,
 $\hat{\mathbf{y}}_{\text{state}}$ via Eqs.\ (\ref{eq:dilconv})---(\ref{eq:heads})
 \State Compute $\mathcal{L}$ via Eq.\ (\ref{eq:loss})
 \State Backpropagate; update $\boldsymbol{\theta}$ with Adam
 \EndFor
 \State Decay $\eta$ with cosine schedule
\EndFor
\State \Return $\boldsymbol{\theta}$
\end{algorithmic}
\end{algorithm}

\subsection{Parameter Budget}\label{sec:arch:params}
% TODO: Table or inline breakdown of 12,052 parameters by component.

%% ============================================================
\section{Experimental Setup}\label{sec:exp}
%% ============================================================

\subsection{Datasets}\label{sec:exp:data}
% TODO: Describe C-MAPSS, NASA Battery, CWRU.
% Cite: Saxena2008, Cheng2020

\begin{table}[!h]
\caption{Dataset summary.}\label{tab:datasets}
\begin{tabular}{lllll}
\toprule
Dataset & Subsets & Sensors & Task & Metric \\
\midrule
NASA C-MAPSS~\cite{Saxena2008} & FD001--FD004 & 14 & RUL regression & RMSE, NASAScore \\
NASA Battery & --- & 4 & SOH / RUL & RMSE, MAE \\
CWRU Bearing & 10 classes & 1 & Fault classification & Accuracy \\
\bottomrule
\end{tabular}
\end{table}

\subsection{Preprocessing}\label{sec:exp:preprocessing}
% TODO: K-Means normalisation for FD002/FD004, z-score for FD001/FD003,
% RUL cap at 125 cycles, window length 64.

\subsection{Baselines}\label{sec:exp:baselines}
% TODO: Describe Ridge Regression, CNN-32 (30.6 KB), LSTM-64 (210.3 KB).

\subsection{Evaluation Metrics}\label{sec:exp:metrics}
% TODO: RMSE, MAE, NASAScore (Eq.~\ref{eq:nasa}), Accuracy,
% Critical-zone RMSE (RUL <= 30).

\subsection{Implementation Details}\label{sec:exp:impl}
% TODO: PyTorch 2.x, Python 3.12, hardware, seeds, epochs.

%% ============================================================
\section{Results and Analysis}\label{sec:results}
%% ============================================================

\subsection{C-MAPSS RUL Estimation}\label{sec:results:cmapss}
% TODO: Discuss Table~\ref{tab:cmapss}.

\begin{table}[!h]
\caption{C-MAPSS RUL estimation (RMSE $\downarrow$, lower is better).%
\tablenotetext{a}{Model size in FP32 bytes.}}
\label{tab:cmapss}
\begin{tabular}{llllll}
\toprule
Model & Size\textsuperscript{a} & FD001 & FD002 & FD003 & FD004 \\
\midrule
Ridge Regression & --- & 44.14 & 54.44 & 42.43 & 55.07 \\
CNN-32 & 30.6\,KB & 26.47 & 40.84 & 24.79 & 39.62 \\
LSTM-64 & 210.3\,KB & 16.28 & 27.95 & 16.02 & \textbf{27.15} \\
\textbf{NanoSentry} & \textbf{47.1\,KB} & \textbf{14.75} & \textbf{26.74} & \textbf{13.95} & 27.67 \\
\bottomrule
\end{tabular}
\end{table}

\subsection{Critical-Zone Analysis}\label{sec:results:critical}
% TODO: Explain the critical-zone metric and discuss Table~\ref{tab:critical}.
% Highlight CNN-32 pathological failure.

\begin{table}[!h]
\caption{Critical-zone RMSE at RUL $\leq$ 30 cycles ($\downarrow$).%
\tablenotetext{a}{CNN-32 critical-zone RMSE exceeds its overall RMSE on FD001
(29.05 vs.\ 26.47), indicating systematic over-estimation near end-of-life.}}
\label{tab:critical}
\begin{tabular}{lllll}
\toprule
Model & FD001 & FD002 & FD003 & FD004 \\
\midrule
\textbf{NanoSentry} & \textbf{4.63} & 5.52 & 3.14 & 6.62 \\
LSTM-64 & 5.27 & \textbf{4.42} & \textbf{2.66} & \textbf{5.94} \\
CNN-32 & 29.05\textsuperscript{a} & 20.67 & 21.50 & 28.18 \\
\bottomrule
\end{tabular}
\end{table}

\subsection{Multi-Domain Results}\label{sec:results:multi}
% TODO: Discuss Table~\ref{tab:multidomain}.

\begin{table}[!h]
\caption{Multi-domain results.}\label{tab:multidomain}
\begin{tabular}{lllll}
\toprule
Dataset & Task & Metric & NanoSentry \\
\midrule
NASA Battery & RUL regression & RMSE & \textbf{1.87} \\
NASA Battery & RUL regression & MAE & \textbf{0.42} \\
CWRU Bearing & Fault classification & Accuracy (\%) & \textbf{99.96} \\
\bottomrule
\end{tabular}
\end{table}

\subsection{NASA Score}\label{sec:results:nasa}
% TODO: Discuss NASAScore table.

\begin{table}[!h]
\caption{NASA asymmetric score ($\downarrow$, lower is better).}\label{tab:nasa}
\begin{tabular}{ll}
\toprule
Subset & NASAScore \\
\midrule
FD001 & 350.47 \\
FD002 & 7405.94 \\
FD003 & 315.86 \\
FD004 & 8677.78 \\
\bottomrule
\end{tabular}
\end{table}

\subsection{Ablation Study}\label{sec:results:ablation}
% TODO: Discuss Table~\ref{tab:ablation}.
% Emphasise CrossSensorGate as dominant component.

\begin{table}[!h]
\caption{Ablation study (RMSE $\downarrow$). Baseline is the full NanoSentry
model trained from the same checkpoint used in Table~\ref{tab:cmapss}.}
\label{tab:ablation}
\begin{tabular}{lllll}
\toprule
Variant & FD001 & $\Delta$ FD001 & FD003 & $\Delta$ FD003 \\
\midrule
\textbf{Full model} & \textbf{14.75} & --- & \textbf{13.95} & --- \\
w/o CrossSensorGate & 17.31 & $+$2.56 & 16.58 & $+$2.63 \\
w/o Skip connection & 15.94 & $+$1.19 & 14.87 & $+$0.92 \\
w/o Dilation & 16.28 & $+$1.53 & 15.41 & $+$1.46 \\
\bottomrule
\end{tabular}
\end{table}

\subsection{Transfer Learning}\label{sec:results:transfer}
% TODO: Discuss Table~\ref{tab:transfer}.
% Frame the negative result as a finding on domain-specificity.

\begin{table}[!h]
\caption{Transfer learning: FD001 pretrain $\rightarrow$ target domain.}
\label{tab:transfer}
\begin{tabular}{llllll}
\toprule
Target & 10\% data & 25\% & 50\% & 100\% & Scratch \\
\midrule
FD003 (RMSE) & 15.08 & 15.35 & 14.64 & 15.31 & \textbf{13.95} \\
Battery (RMSE) & 3.84 & 3.83 & 2.48 & 3.75 & \textbf{1.87} \\
\bottomrule
\end{tabular}
\end{table}

%% --- FIGURE 2: RUL prediction curves ---
\begin{figure}[!t]
 \centering
 % \includegraphics[width=\columnwidth]{figs/rul_curves.pdf}
 \fbox{\parbox{0.9\columnwidth}{\centering\vspace{2cm}
 [PLACEHOLDER: RUL prediction curves --- NanoSentry vs LSTM-64
 on one FD001 test engine. Predicted vs.\ true RUL over cycles.]\vspace{2cm}}}
 \caption{Predicted versus true RUL for a representative FD001 test engine.
 NanoSentry tracks the degradation trajectory with lower end-of-life
 deviation than LSTM-64 despite being 4.5$\times$ smaller.}
 \label{fig:rul_curves}
\end{figure}

%% --- FIGURE 3: Critical-zone scatter ---
\begin{figure}[!t]
 \centering
 % \includegraphics[width=\columnwidth]{figs/critical_zone.pdf}
 \fbox{\parbox{0.9\columnwidth}{\centering\vspace{2cm}
 [PLACEHOLDER: Critical-zone scatter plot. All three models,
 actual vs.\ predicted RUL for RUL $\leq$ 30 cycles, FD001.]\vspace{2cm}}}
 \caption{Critical-zone scatter (RUL $\leq$ 30 cycles, FD001). Each point
 is one test-set sample. CNN-32 (red) exhibits severe over-estimation bias;
 NanoSentry (blue) clusters tightly around the identity line.}
 \label{fig:critical}
\end{figure}

%% ============================================================
\section{Discussion}\label{sec:discussion}
%% ============================================================

\subsection{Accuracy--Size Trade-Off}\label{sec:disc:tradeoff}
% TODO: Discuss 4.5x size advantage with competitive RMSE.
% Contextualise against MCU deployment constraints.

\subsection{Critical-Zone Behaviour}\label{sec:disc:critical}
% TODO: Explain why CNN-32 pathologically fails in the critical zone.
% Discuss safety implications for maintenance scheduling.

\subsection{Transfer Learning and Domain Specificity}\label{sec:disc:transfer}
% TODO: Argue that the negative transfer result is itself a contribution.
% Propose federated PHM as a future direction.
% Cite: Almeida2026 (FL-Offshore paper)

%% ============================================================
\section{Conclusions}\label{sec:conclusions}
%% ============================================================

% TODO: Three short paragraphs:
% 1. What was proposed and demonstrated.
% 2. Three key findings (one sentence each).
% 3. Future work: uncertainty quantification, federated learning,
% real-time embedded validation.

%% ============================================================
%% BIBLIOGRAPHY
%% ============================================================

%% Option A (recommended): BibTeX
%\bibliographystyle{wileyNJD-AMA}
%\bibliography{references}

%% Option B: Inline (paste your .bbl output here if needed)
\begin{thebibliography}{99}

\bibitem{Lei2018}
Lei Y, Li N, Guo L, Li N, Yan T, Lin J.
Machinery health prognostics: a systematic review from data acquisition
to RUL prediction.
\textit{Mech Syst Signal Process.} 2018;104:799--834.

\bibitem{Saxena2008}
Saxena A, Goebel K, Simon D, Eklund N.
Damage propagation modeling for aircraft engine run-to-failure simulation.
In: \textit{2008 International Conference on Prognostics and Health
Management}; 2008:1--9.

\bibitem{Zhang2019}
Zhang W, Yang D, Wang H.
Data-driven methods for predictive maintenance of industrial equipment:
a survey.
\textit{IEEE Syst J.} 2019;13(3):2213--2227.

\bibitem{Zheng2017}
Zheng S, Ristovski K, Farahat A, Gupta C.
Long short-term memory network for remaining useful life estimation.
In: \textit{2017 IEEE International Conference on Prognostics and Health
Management (ICPHM)}; 2017:88--95.

\bibitem{Mo2021}
Mo Y, Wu Q, Li X, Huang B.
Remaining useful life estimation via transformer encoder enhanced by
a gated convolutional unit.
\textit{J Intell Manuf.} 2021;32(7):1997--2006.

\bibitem{Elsheikh2019}
Elsheikh A, Yacout S, Ouali M-S.
Bidirectional handshaking LSTM for remaining useful life prediction.
\textit{Neurocomputing.} 2019;323:148--156.

\bibitem{Li2020}
Li H, Zhao W, Zhang Y, Zio E.
Remaining useful life prediction using multi-scale deep convolutional
neural network.
\textit{Appl Soft Comput.} 2020;89:106113.

\bibitem{Zhang2023atcn}
Zhang Q, Liu Q, Ye Q.
An attention-based temporal convolutional network method for predicting
remaining useful life of aero-engine.
\textit{Eng Appl Artif Intell.} 2024;127:107241.

\bibitem{Zhang2023imdssn}
Zhang J, Li X, Tian J, Luo H, Yin S.
An integrated multi-head dual sparse self-attention network for
remaining useful life prediction.
\textit{Reliab Eng Syst Saf.} 2023;233:109096.

\bibitem{Lin2023}
Lin R, Wang H, Xiong M, Hou Z, Che C.
Attention-based gate recurrent unit for remaining useful life prediction
in prognostics.
\textit{Appl Soft Comput.} 2023;143:110419.

\bibitem{Liu2021}
Liu H, Liu Z, Jia W, Lin X.
Remaining useful life prediction using a novel feature-attention-based
end-to-end approach.
\textit{IEEE Trans Ind Inform.} 2021;17(2):1197--1207.

\bibitem{Chen2024}
Chen X.
A novel transformer-based DL model enhanced by position-sensitive
attention and gated hierarchical LSTM for aero-engine RUL prediction.
\textit{Sci Rep.} 2024;14:10061.

\bibitem{Gong2022}
Gong R, Li J, Wang C.
Remaining useful life prediction based on multisensor fusion and
attention TCN-BiGRU model.
\textit{IEEE Sens J.} 2022;22(21):21101--21110.

\bibitem{Wang2024}
Wang F, Zhai Z, Zhao Z, Di Y, Chen X.
Physics-informed neural network for lithium-ion battery degradation
stable modeling and prognosis.
\textit{Nat Commun.} 2024;15:4332.

\bibitem{Gong2024}
Gong L, Zhang Z, Li X, et al.
State of health estimation of lithium-ion batteries based on machine
learning with mechanical-electrical features.
\textit{Batteries Supercaps.} 2024;7(10):e202400201.

\bibitem{Richardson2017}
Richardson RR, Osborne MA, Howey DA.
Gaussian process regression for forecasting battery state of health.
\textit{J Power Sources.} 2017;357:209--219.

\bibitem{Cheng2020}
Cheng C, Ma G, Zhang Y, et al.
A deep learning-based remaining useful life prediction approach
for bearings.
\textit{IEEE/ASME Trans Mechatron.} 2020;25(3):1243--1254.

\bibitem{Han2021}
Han T, Li Y-F, Qian M.
A hybrid generalization network for intelligent fault diagnosis of
rotating machinery under unseen working conditions.
\textit{IEEE Trans Instrum Meas.} 2021;70:1--11.

\bibitem{Zhao2020}
Zhao M, Zhong S, Fu X, Tang B, Pecht M.
Deep residual shrinkage networks for fault diagnosis.
\textit{IEEE Trans Ind Inform.} 2020;16(7):4681--4690.

\bibitem{Wen2018}
Wen L, Li X, Gao L, Zhang Y.
A new convolutional neural network-based data-driven fault diagnosis method.
\textit{IEEE Trans Ind Electron.} 2018;65(7):5990--5998.

\bibitem{Zhang2018}
Zhang W, Li C, Peng G, Chen Y, Zhang Z.
A deep convolutional neural network with new training methods for bearing
fault diagnosis under noisy environment and different working load.
\textit{Mech Syst Signal Process.} 2018;100:439--453.

\bibitem{Liao2023}
Liao J-X, Dong H-C, Sun Z-Q, et al.
Attention-embedded quadratic network (Qttention) for effective and
interpretable bearing fault diagnosis.
\textit{IEEE Trans Instrum Meas.} 2023;72:1--13.

\bibitem{Liang2022}
Liang H, Cao J, Zhao X.
Multi-scale dynamic adaptive residual network for fault diagnosis.
\textit{Measurement.} 2022;188:110397.

\bibitem{Zhou2025}
Zhou L, Wang H.
Multi-task model of adaptive multi-scale feature fusion and adaptive
mixture-of-experts for equipment remaining useful life prediction and
fault diagnosis.
\textit{Expert Syst Appl.} 2025;272:126807.

\bibitem{Cheng2021}
Cheng H, Kong X, Chen G, Wang Q, Wang R.
Transferable convolutional neural network based remaining useful life
prediction of bearing under multiple failure behaviors.
\textit{Measurement.} 2021;168:108286.

\bibitem{Li2021}
Li X, Zhang W.
Deep learning-based partial domain adaptation method on intelligent
machinery fault diagnostics.
\textit{IEEE Trans Ind Electron.} 2021;68(5):4351--4361.

\bibitem{Ding2023}
Ding W, Li J, Mao W, Meng Z, Shen Z.
Rolling bearing remaining useful life prediction based on dilated
causal convolutional DenseNet and an exponential model.
\textit{Reliab Eng Syst Saf.} 2023;232:109072.

\bibitem{Guo2017}
Guo L, Li N, Jia F, Lei Y, Lin J.
A recurrent neural network based health indicator for remaining useful
life prediction of bearings.
\textit{Neurocomputing.} 2017;240:98--109.

\bibitem{Ochella2024}
Ochella S, Dinmohammadi F, Shafiee M.
Bayesian neural networks for uncertainty quantification in remaining
useful life prediction.
\textit{Adv Mech Eng.} 2024;16(7):16878132241239802.

\bibitem{Jiang2026}
Jiang Z, Wei Y, Li X.
Remaining useful life prediction across different operational conditions
based on domain adaptation feature transfer.
\textit{Sci Rep.} 2026. doi:10.1038/s41598-026-49358-6.

\end{thebibliography}

\end{document}
